% !TEX root = ../main.tex
\section{Background and theory}

\hypertarget{dynamic-inconsistency-and-bellmans-principle}{%
\subsection{Dynamic inconsistency and Bellman's principle}\label{dynamic-inconsistency-and-bellmans-principle}}

Consider a planner who at time zero chooses a sequence of decisions---here, harvest levels by stratum and period---over a finite horizon of \(T\) periods, so as to maximize an objective (discounted net present value, NPV) subject to constraints, including a harvest-flow constraint that links consecutive periods. The open-loop solution is a complete plan: a decision for every period, chosen once, at time zero, from the initial forest state.

Now suppose the plan is executed one period at a time by a sequence of planners, each inheriting the forest state realized by her predecessors and each solving the \emph{same} optimization problem from that state over the remaining horizon. Bellman's principle of optimality states that an optimal policy has the property that, whatever the initial state and initial decision, the remaining decisions must constitute an optimal policy with regard to the state resulting from the first decision \citep{bellman1957dynamic}. A plan that satisfies this property is \emph{dynamically consistent}: its tail is optimal for every subproblem it passes through, so each successive planner, re-solving from the realized state, voluntarily continues the plan. A plan that fails the property is \emph{dynamically inconsistent}: at some point the re-solving planner finds the plan's tail suboptimal for the state she has inherited and deviates from it. The open-loop projection of period-by-period harvests is then not a forecast of anything that will happen; it is a statement about a counterfactual world in which precommitment is possible.

Formally, let \(x = (x_1, \ldots, x_T)\) be a sequence of decisions and let the open-loop problem from state \(s_\tau\) over the remaining horizon be \(V_\tau(s_\tau) = \max_{x_\tau, \ldots, x_T} \sum_{t=\tau}^{T} \delta^{\,t-\tau} \, u_t(x_t)\) subject to the model's constraints, with \(\delta\) the per-period discount factor. The open-loop plan announced at time zero is \(x^{0} = (x_1^{0}, \ldots, x_T^{0}) = \arg\max V_1(s_1)\). Under sequential replanning, the planner at period \(\tau\) observes the realized state \(s_\tau = f(s_{\tau-1}, x_{\tau-1})\) (the forest left by her predecessor's decision) and re-solves \(V_\tau(s_\tau)\). The plan \(x^{0}\) is dynamically consistent if its tail is optimal for every subproblem it passes through, i.e., if \((x_\tau^{0}, \ldots, x_T^{0})\) remains a maximizer of \(V_\tau(s_\tau)\) for all \(\tau\); otherwise some planner deviates and the plan is not followed.

The practical stakes are not subtle. Long-term forest plans are used to set allowable cuts, to justify regulation of harvest flow, and to price constraints and land uses through shadow prices. If the plan is dynamically inconsistent, the projected harvest levels will not be delivered, the projected residual forest will not exist, and the shadow prices---which are properties of the particular optimal solution that will not be followed---are biased as measures of trade-offs along any trajectory that will actually occur.

\hypertarget{the-open-loop-lp-forest-planning-problem}{%
\subsection{The open-loop LP forest-planning problem}\label{the-open-loop-lp-forest-planning-problem}}

The standard strategic timber harvest-scheduling model is strata-based: the forest is partitioned into development types (strata) defined by attributes such as site, species composition, age, and management prescription, and the decision variables allocate stratum area to management activities over a horizon of discrete periods. The classical formulations are those of \citet{johnson1977techniques}. In \emph{Model I}, each existing stratum is assigned a complete management schedule as a distinct activity, so that the trajectory of every initial cohort is represented explicitly. In \emph{Model II}, decisions are expressed more compactly as area flows into a smaller set of prescription activities; \emph{Model III} aggregates further still, into a network flow over age classes. We refer the reader to \citep[ch.~16]{gunn2007models} for a careful statement of the Model I/II/III distinction and its lineage.

A word on model form and size is warranted, because a misconception about it circulates. Model II is widely described as the more compact formulation on the strength of its smaller variable count (Johnson and Scheurman 1977), and it is sometimes asserted that Model II reduces the size and difficulty of harvest-scheduling LPs relative to Model I. That expectation does not survive contact with modern instances. In the first published comparison of the two formulations on realistic spatial problems, \citet{martin2017comparing} showed that Model II's apparent compactness advantage evaporates---and reverses---once the problem carries realistic spatial resolution or many strata: Model II's per-strata constraint network makes its matrix denser and its solution time grow rapidly with added structure, while Model I's cost grows only marginally. On their full case study, Model I reached comparable optimal solutions in roughly a tenth of Model II's solution time. The authors' conclusion is blunt: the common expectation that Model II is the smaller, easier model is \emph{refuted} for all but the smallest or most weakly structured problems. This bears directly on the reproduction's formulation choice (below).

Two ingredients of the standard formulation matter for what follows. First, the objective is discounted NPV: timber values and costs in each period are weighted by a discount factor, so early harvest is preferred to late harvest, all else equal. Second, harvest-flow constraints---typically even-flow or bounded-deviation constraints on total volume between consecutive periods---create the only direct link between decisions in different periods. The economics of such harvest-flow constraints---the even-flow constraint and the allowable-cut effect it enables---have a long critical literature \citep{thompson1966traditional,schweitzer1972allowable,luckert1995allowable,hegan2000economic}. Without such a link, the multi-period problem decomposes and open-loop and sequential solutions coincide. With it, the current period's decision changes the feasible set---and the shadow prices---of every future subproblem.

Writing \(H_t\) for the total harvest volume in period \(t\) and suppressing the per-stratum detail, the open-loop problem maximizes discounted net revenue subject to a harvest-flow constraint linking consecutive periods:
\begin{equation}
\max_{H_1, \ldots, H_T \geq 0} \;\; \sum_{t=1}^{T} \delta^{\,t} \, r_t \, H_t
\qquad \text{s.t.} \qquad
(1 - \alpha)\, H_t \;\le\; H_{t+1} \;\le\; (1 + \beta)\, H_t, \;\; t = 1, \ldots, T-1,
\label{eq:openloop}
\end{equation}
where \(r_t\) is net revenue per unit volume in period \(t\) and \(\delta = (1+i)^{-1}\). The harvest-flow constraint is the classic FORPLAN bounded-deviation ("sequential flow") form \citep[eq.~3-4, 3-5]{daugherty1991credibility}: \(\alpha\) is the maximum allowed fractional \emph{decrease} and \(\beta\) the maximum allowed fractional \emph{increase} between consecutive periods. Its special cases span the standard harvest-flow policies: non-declining yield (\(\alpha = 0\), no increase bound; the yield may not fall), bounded decline (\(\alpha > 0\), no increase bound), and symmetric bounded deviation (\(\alpha = \beta > 0\)). It is this between-period link---and only this link---that prevents the multi-period problem from decomposing into independent single-period problems, and so it is the structural feature on which the consistency question turns.

\citet{daugherty1991credibility} conducted his analysis using a Model II formulation. The reproduction reported here builds the case study as a Model I model in \texttt{ws3}. The distinction matters for computational detail but not for the consistency question at issue: dynamic inconsistency is a property of the interaction among the objective, the flow constraint, and the initial forest structure, not of the variable-aggregation scheme. Moreover, the Model I choice is not merely a convenience: on the evidence of \citet{martin2017comparing}, Model I is at least as computationally efficient as Model II on realistic, well-structured problems, so nothing is sacrificed computationally by reproducing the analysis in the stand-based form. We nonetheless note the formulation difference explicitly as a fidelity caveat rather than a deviation of substance.

\hypertarget{daughertys-thesis}{%
\subsection{Daugherty's thesis}\label{daughertys-thesis}}

\citet{daugherty1991credibility} posed the credibility question directly: are the plans produced by LP-based forest-planning models, solved open-loop, plans that a sequence of future planners operating the same model would actually follow? The thesis (i) formalized dynamic consistency for LP forest-planning models, (ii) derived conditions under which open-loop optimal plans fail the consistency requirements, and (iii) demonstrated the phenomenon numerically by iterative LP simulation of sequential replanning on a real case study---the Umpqua National Forest FORPLAN planning model documented in \citep{usda1987umpqua}---across a designed experiment varying the initial forest condition, the harvest policy, and the interest rate. The thesis found dynamic inconsistency to be pervasive rather than exceptional, and traced it to the combination of the harvest-flow constraint with a disequilibrium initial forest containing financially over-mature and negatively valued strata. \citet[pp.~331--332]{gunn2007models}, citing \citet{daugherty1991credibility}, summarizes the practical symptom memorably as the ``declining non-declining yield'': a plan that promises a non-declining harvest trajectory while the sequential re-optimization it invites systematically revises the attainable harvest downward.
